\section{Conformance requirements}
\label{sec:conformance}

The implementation is built against these requirements. Each is
cheap to meet while the system is small, and expensive to retrofit once runs, reports and
agreements exist.

\begin{description}[leftmargin=2.6em,itemsep=0.5em]
  \item[C1: no patient-level entity.] The system stores no identifier, pseudonym, row or key at
        patient granularity. Demand appears only as aggregates per facility, commodity and
        period.
  \item[C2: an audit record per action.] Starting a run, changing a configuration, reading a
        model output and exporting a report each produce a record naming the operator, the object,
        the time and the configuration hash.
  \item[C3: deletion by run.] A run's raw inputs, model versions, evaluation snapshots and
        audit records must be deletable as one operation, and the deletion must be
        demonstrable.
  \item[C4: one purpose per dataset version.] Each dataset version states the purpose it was
        collected for. A second purpose needs a new agreement.
  \item[C5: synthetic until agreed.] Until a data-sharing agreement and ethics clearance
        exist, every artifact carries the word synthetic, and the label travels inside the
        file.
  \item[C6: one name per object.] Every policy, entity and parameter keeps a single name wherever
        it appears: in this document, in the configuration of a run, and in any report produced
        from it. The implementation checks its own vocabulary against the one defined here and
        reports a name that matches nothing.
\end{description}

\noindent
C1, C5 and C6 bind the first build. The others are what a deployment would need, listed now
so that the prototype does not have to be dismantled first.
